\chapter{Introduction}
\label{ch:introduction}

\section{Problem Statement}
\label{sec:intro-problem}

A digital VLSI circuit must satisfy two independent kinds of correctness.
The first is \emph{logical} correctness: for every legal input combination,
the circuit must eventually settle to the specified output values. The
second, and the subject of this report, is \emph{timing} correctness: every
signal transition must arrive at its destination early enough to be
correctly captured, and no path in the circuit may take longer than the
timing budget allotted to it. A circuit that is logically correct but
functionally simulated only at a handful of input vectors can still fail in
silicon because some rarely-exercised path violates its timing constraint;
timing correctness is therefore not amenable to verification by input-vector
simulation, since the number of paths through a circuit of even modest size
is combinatorially large and no practical vector set can be guaranteed to
sensitize the worst one.

\emph{Static timing analysis} (STA) sidesteps this problem entirely. Rather
than simulating logic values, STA treats the circuit as a directed acyclic
graph (DAG) of timing arcs and propagates \emph{worst-case} delay through
that graph, independent of the actual logic values that would appear on any
particular clock cycle. Every combinational path is implicitly covered by
construction, because the propagation is a graph traversal over every arc,
not an enumeration of input vectors. The output of STA is, for every node in
the circuit, an \emph{arrival time} (the latest moment a signal transition
can occur at that node) and a \emph{required time} (the latest moment it is
permitted to occur without violating a downstream constraint); their
difference is \emph{slack}, and a negative slack anywhere in the circuit is
a timing violation.

Once a violation is identified, one of the standard ways to close it without
changing the circuit's logic is \emph{gate sizing}: replacing a violating
path's gates with larger-drive-strength variants from the same logic family
in the standard-cell library. A larger gate has lower output resistance and
therefore a lower propagation delay for the same load, at the cost of a
larger input capacitance (which increases the load, and hence the delay, of
whatever gate drives it), a larger area, and higher leakage and dynamic
power. Because these effects compose across the whole circuit, sizing one
gate can shift the critical path elsewhere or worsen a previously
non-critical path, and a purely trial-and-error sizing process is both slow
and difficult to reason about. \emph{Logical effort} is the theoretical
framework used in this project to guide that search: it decomposes a path's
delay into a logic-topology-dependent term (how many stages, and how
electrically ``expensive'' each stage's logic function is) and an
electrical-loading term (how much capacitance each stage must drive relative
to its own size), which together identify \emph{which} gate on a violating
path is most under-sized for its load and therefore the best candidate for
resizing.

\section{Project Scope and Deliverables}
\label{sec:intro-scope}

The assignment specification requires a software tool that, given a gate-level
netlist, a sized standard-cell library, and a configuration file, must:

\begin{enumerate}
    \item parse the inputs and validate the netlist's structural
          correctness, reporting every detected error class distinctly;
    \item build a combinational DAG and topologically order it;
    \item compute output load capacitance and propagation delay per gate,
          separately for rising and falling output transitions;
    \item perform forward (arrival-time) and backward (required-time)
          timing propagation, correctly handling positive-unate,
          negative-unate, and non-unate timing arcs;
    \item report per-node slack, circuit-level worst negative slack (WNS)
          and total negative slack (TNS), and the top-$K$ least-slack
          critical paths;
    \item perform a logical-effort decomposition of each critical path and
          use it to guide a search for an improved gate-sizing assignment,
          honoring area and power budgets;
    \item estimate power and area before and after optimization;
    \item compare the logical-effort-guided search against a baseline
          sizing method that does not use logical effort; and
    \item (bonus) perform a Monte Carlo statistical timing analysis that
          models manufacturing process variation.
\end{enumerate}

\section{Contributions Beyond the Base Assignment}
\label{sec:intro-contributions}

While the base requirements above were implemented in full, two additional
components were developed that go beyond what the assignment strictly
requires, and this report documents them at the same level of detail as the
required components:

\begin{itemize}
    \item \textbf{A synthetic benchmark generation and evaluation
          framework} (\cref{ch:benchmarking}), which programmatically
          constructs large numbers of controlled, repairable timing-violation
          scenarios --- including a certified best-known reference solution
          obtained via multi-start beam search --- and evaluates every
          sizing heuristic's decisions against an independent counterfactual
          oracle, reporting statistically rigorous metrics (Wilson
          confidence intervals on repair rate, decision regret, and
          stratified breakdowns by circuit size, depth, fan-out, and
          reconvergence). This addresses a limitation of evaluating sizing
          heuristics only on the fifteen fixed circuits supplied with the
          assignment: fifteen circuits are too few to distinguish a
          genuinely better heuristic from one that simply got lucky on a
          particular benchmark.
    \item \textbf{An interactive, browser-based topology viewer}
          (\cref{ch:viewer}), which renders a circuit's structure with
          pan/zoom navigation, slack-based coloring, critical-path
          highlighting, and a toggle between the pre-optimization and
          post-optimization gate-sizing state of the same circuit.
\end{itemize}

These additions are clearly delineated from the required deliverables
throughout this report so that grading against the assignment specification
remains straightforward.

\section{Report Roadmap}
\label{sec:intro-roadmap}

\Cref{ch:background} reviews the theoretical background needed to read the
rest of the report: the STA graph model, the RC delay model, logical effort,
power/area estimation, and the statistical variation model used for Monte
Carlo analysis. \Cref{ch:architecture} describes the software architecture
and the overall data flow, including the figure referenced throughout the
rest of the report. \Cref{ch:validation,ch:sta,ch:logical-effort,ch:power-area}
document, in the order they execute, the input pipeline, the STA engine, the
logical-effort analysis, and the power/area estimator.
\Cref{ch:optimization} documents the cost function, the four sizing
heuristics, and their acceptance/termination rules. \Cref{ch:monte-carlo}
documents the statistical timing analysis. \Cref{ch:benchmarking,ch:viewer}
document the two extensions described above. \Cref{ch:results} presents
results on both the fixed and synthetic benchmark suites and \cref{ch:conclusion} concludes.
